\section{中央处理器与图形处理器的设计取向}

\subsection{先区分延迟与吞吐量}

比较中央处理器和图形处理器之前，必须先区分两个性能指标。

\begin{itemize}
  \item \term{延迟}{latency}描述单个任务从开始到完成所经历的时间。若任务在时刻 $t_{\mathrm{start}}$ 开始、在 $t_{\mathrm{finish}}$ 完成，则
  \begin{equation*}
    L=t_{\mathrm{finish}}-t_{\mathrm{start}}.
  \end{equation*}
  延迟越低，单个请求越快得到结果。
  \item \term{吞吐量}{throughput}描述单位时间内完成的任务数量。若时间区间 $\Delta t$ 内完成 $N$ 个任务，则平均吞吐量为
  \begin{equation*}
    R=\frac{N}{\Delta t}.
  \end{equation*}
  吞吐量越高，系统在一段时间内处理的总工作越多。
\end{itemize}

低延迟与高吞吐量并不等价。一个流水线可以让单个任务经过多个阶段，因此单任务延迟较长；当流水线填满后，各阶段同时处理不同任务，系统仍可在每个周期完成一个或多个结果，从而获得高吞吐量。

\begin{modelbox}
设一条理想流水线包含 $d$ 个阶段，时钟频率为 $f$。单个任务的流水线延迟约为 $d/f$，但流水线达到稳态后，每条执行通道的理想结果产生速率可以接近 $f$。增加流水线深度可能不降低单任务延迟，却能支持更高的稳态工作速率，从而形成图形处理器“长延迟、高吞吐量”的组合。
\end{modelbox}

\subsection{中央处理器：面向低延迟的设计}

中央处理器采用\term{面向延迟的设计}{latency-oriented design}，主要优化少量指令流的快速完成，其芯片面积通常分配给以下结构。

\subsubsection{少量但功能强的算术逻辑单元}

\term{算术逻辑单元}{Arithmetic Logic Unit, ALU}执行整数、逻辑及部分数值运算。中央处理器核心包含较少但较复杂的执行单元，目标是降低单条指令或短指令序列的完成时间。复杂的乱序执行与多种执行端口也属于同一低延迟取向。

\subsubsection{较大的缓存层次}

主存访问延迟远高于算术操作。中央处理器投入大量芯片面积构建多级\term{缓存}{cache}，把频繁访问的数据保留在靠近执行单元的位置。缓存命中时，原本较长的主存访问被转换为较短的缓存访问，从而降低单线程停顿。

\subsubsection{复杂控制逻辑}

中央处理器使用复杂控制结构减少依赖与分支带来的等待：
\begin{itemize}
  \item \term{分支预测}{branch prediction}提前猜测控制流方向，尽量避免处理流水线因等待分支结果而空闲；
  \item \term{数据转发}{data forwarding}把某条指令刚产生的结果直接送给后续相关指令，减少等待数据写回再读取的时间；
  \item 适度多线程用于隐藏相对较短的延迟，但线程数量通常远少于图形处理器。
\end{itemize}

中央处理器还通常采用较高时钟频率。综合这些机制，它非常适合控制密集、分支复杂、串行依赖明显或任务规模较小的程序部分。

\subsection{图形处理器：面向高吞吐量的设计}

图形处理器采用\term{面向吞吐量的设计}{throughput-oriented design}，主要优化大量相似任务的总完成速率。

\subsubsection{大量较小的算术单元}

图形处理器把更多芯片面积用于计算通道，因而能同时保持大量算术操作在执行流水线中。单个操作或单个线程的延迟可能高于中央处理器，但当工作规模足够大时，所有线程的总完成时间可以显著缩短。

\subsubsection{深流水线与海量线程}

图形处理器通过深流水线提高稳态吞吐量，并维护大量可执行线程。当一组线程等待长延迟事件，例如全局内存数据返回时，调度器可以选择另一组已经具备操作数的线程执行。

这个机制称为\term{延迟隐藏}{latency hiding}。它没有缩短那次内存访问本身的物理延迟；它通过让其他工作占用原本会空闲的执行周期，减少延迟对总体吞吐量的影响。若程序没有足够多的独立线程，或所有线程同时等待同一瓶颈，延迟就无法被有效隐藏。

\subsubsection{相对较小的缓存与较简单的控制}

图形处理器也包含缓存，但相对于中央处理器，它把更大比例的面积用于算术单元、寄存器与线程执行资源。控制结构也偏向规则执行，因为最有效的工作负载通常让许多线程对不同数据执行相同或相近的指令序列。其时钟频率通常适中，性能主要来自并行宽度与总吞吐量，而不是极高频率。

\begin{figure}[H]
  \centering
  \includegraphics[width=0.98\textwidth]{cpu_gpu_design.png}
  \caption{中央处理器的低延迟取向与图形处理器的高吞吐量取向}
  \label{l1:fig:cpu-gpu-design}
\end{figure}

\subsection{芯片面积分配反映了不同目标}

图~\ref{l1:fig:cpu-gpu-design} 的方块图不是精确的晶体管布局，而是面积预算的概念图。中央处理器把大量面积用于缓存和复杂控制，以减少少量线程的等待；图形处理器把更多面积复制为大量计算通道，以提高规则并行工作负载的总处理速率。

这种权衡可以从机会成本理解。若增加缓存和控制逻辑，就会减少可放置的算术通道；若增加算术通道，就必须接受较少的单线程优化资源。两类处理器都不是绝对“更快”，它们针对不同目标函数优化。

两张芯片照片可以进一步说明粒度差异：第十代英特尔酷睿处理器 (10th Gen Intel Core processor) 在 14 nm 工艺下包含 10 个中央处理器核心；英伟达 GK110 在 28 nm 工艺下标有 2880 个统一计算设备架构核心。较老工艺的 GK110 仍能容纳数量大得多的“核心”，正是因为这里的统一计算设备架构核心远小于完整中央处理器核心。该对比用于说明结构粒度，不能仅按数字推断两张芯片的通用性能。

\begin{table}[H]
\centering
\caption{中央处理器与图形处理器的高层对比}
\begin{tabularx}{\textwidth}{L{0.22\textwidth}Y Y}
\toprule
\textbf{维度} & \textbf{中央处理器} & \textbf{图形处理器} \\
\midrule
主要目标 & 降低单个线程或请求的延迟 & 提高大批相似任务的总吞吐量 \\
算术资源 & 少量、复杂、功能强的执行单元 & 大量、较小、深流水化的执行单元 \\
缓存 & 较大，用于降低平均访存延迟 & 相对较小，把更多面积留给计算资源 \\
控制 & 分支预测、数据转发等复杂机制 & 偏向规则控制和大规模线程调度 \\
并发方式 & 少量强线程，适度多线程 & 海量线程，通过切换隐藏长延迟 \\
典型优势部分 & 串行依赖、复杂控制、低延迟响应 & 数据并行、规则数值计算、高吞吐批处理 \\
\bottomrule
\end{tabularx}
\end{table}

\subsection{“核心”一词在两类处理器中并不等价}

\term{核心}{core}在两类处理器中具有不同的结构粒度和语义。

中央处理器核心是相对完整的处理单元，通常包含多个执行单元、寄存器、至少一级私有缓存和复杂控制逻辑。它能够独立推进一个或多个通用指令流。

英伟达所称的\term{统一计算设备架构核心}{CUDA core}则是更细粒度的算术数据通路，可以近似理解为流式处理器中的一个算术执行单元。它不具备中央处理器核心那样完整的缓存与控制结构。因此，“10 个中央处理器核心”与“2880 个统一计算设备架构核心”不能按核心数量直接比较；两者的粒度、能力和组织方式完全不同。

\subsection{高性能系统通常同时使用中央处理器与图形处理器}

异构系统中的任务划分原则是：
\begin{itemize}
  \item 对延迟敏感、难以并行化的串行部分由中央处理器执行；
  \item 对吞吐量敏感、含有大量相互独立工作的并行部分由图形处理器执行。
\end{itemize}

两类处理器在各自擅长的代码上可能相差十倍以上。该数值是概念性数量级说明，不是对任意程序的固定保证。实际加速取决于可并行比例、数据传输、内存访问、分支规则性、问题规模和实现质量。

一个典型异构应用的执行过程是：中央处理器读取与准备数据、决定计算阶段和发起图形处理器任务；图形处理器执行规模最大的并行计算；结果再被中央处理器用于后处理、控制或输出。两种设备之间的协同本身会引入数据移动与同步开销，因此正确的任务划分不仅要比较计算速度，还要判断把任务迁移到加速器是否足以抵消这些额外成本。

\begin{keybox}
中央处理器与图形处理器的根本差异可以概括为：中央处理器用更多硬件资源缩短少数线程的关键路径，图形处理器用更多并行执行资源扩大单位时间内完成的工作总量。图形处理器编程的首要条件不是“有很多算术”，而是“有很多能够同时推进、且执行结构足够规则的算术”。
\end{keybox}
